\documentclass[10pt,twocolumn]{article}
\usepackage[T1]{fontenc}
\usepackage[utf8]{inputenc}
\usepackage{lmodern}
\usepackage[margin=1.8cm,top=2.4cm,bottom=2cm,columnsep=0.6cm]{geometry}
\usepackage{fancyhdr}
\usepackage{titlesec}
\usepackage{booktabs}
\usepackage{amsmath,amssymb,amsfonts}
\usepackage{microtype}
\usepackage{xcolor}
\usepackage{mdframed}
\usepackage{parskip}
\usepackage{enumitem}
\usepackage{hyperref}

\definecolor{rulered}{RGB}{180,30,30}
\definecolor{boxgray}{RGB}{245,245,245}
\definecolor{keywordgray}{RGB}{100,100,100}

\pagestyle{fancy}
\fancyhf{}
\fancyhead[L]{\textsc{\small Claude Mag}}
\fancyhead[C]{\small\textit{Machine Learning Research Digest}}
\fancyhead[R]{\small 2026-10-05}
\fancyfoot[C]{\small\thepage}
\renewcommand{\headrulewidth}{0.8pt}

\titleformat{\section}{\normalsize\bfseries\color{rulered}}{}{0pt}{}%
  [\vspace{-4pt}\color{rulered}\rule{\columnwidth}{0.4pt}]
\titlespacing{\section}{0pt}{8pt}{4pt}

\hypersetup{colorlinks=true,urlcolor=rulered,linkcolor=black}

\begin{document}
\setlength{\parindent}{0pt}
\setlength{\parskip}{4pt}

{\small\color{keywordgray}\textbf{Keywords:}
  mechanistic interpretability \textbullet{}
  non-linear manifolds \textbullet{}
  concept alignment \textbullet{}
  LLM representations}\par\vspace{4pt}

{\LARGE\bfseries\raggedright
  Beyond Linear Concepts: Discovering and Aligning Non-Linear Concept Manifolds
  in Large Language Models}\par\vspace{4pt}

{\large\itshape\raggedright
  Modelling Concepts as Low-Dimensional Riemannian Manifolds Reveals
  Consistent Block Structures and a Syntax-to-Semantics Transition
  Universal Across Architectures}\par\vspace{6pt}

{\small\textbf{By} Tido Specht, Elias Benedict Krey, Nils Neukirch
  \& Nils Strodthoff (University of Oldenburg; DFKI)}\par\vspace{2pt}

{\small\color{keywordgray}arXiv:2610.01821 \textbullet{} October 2026}
\par\vspace{2pt}\noindent{\color{rulered}\rule{\columnwidth}{0.6pt}}\vspace{6pt}

Adapting Non-Linear Multi-Dimensional Concept Discovery (NLMCD) from computer vision
to token-level LLM activations and introducing the Concept-Based Alignment (CBA)
score, this work reveals that concept organization in LLMs exhibits consistent
block structures across layers and architectures, with representations staying
syntax-dominated through 70\% of the network before transitioning to mixed
syntactic-semantic structure in late layers---patterns invisible to linear methods.

\begin{mdframed}[backgroundcolor=boxgray,linecolor=rulered,linewidth=1pt,
    innerleftmargin=6pt,innerrightmargin=6pt,
    innertopmargin=6pt,innerbottommargin=6pt]
\textbf{\small AT A GLANCE}\par\vspace{4pt}
\begin{description}[leftmargin=0pt,labelsep=4pt,itemsep=2pt,parsep=0pt,
    font=\small\bfseries]
  \item[Problem:] Existing MI methods assume linear concept representations,
    missing the non-linear geometric structure of LLM activations.
  \item[Why it matters:] Linear methods may produce incomplete or misleading
    interpretations of how concepts are organized in large models.
  \item[Method:] NLMCD fits low-dimensional manifolds to token activations;
    CBA score (generalized Rand index) compares manifolds across layers and models
    without requiring explicit feature matching.
  \item[Result:] Consistent block structures at intermediate layers; syntax-dominated
    through 70\% of network then mixed syntactic-semantic; CBA is 15--30\% more
    sensitive than PCA or CKA.
  \item[Evidence:] Five LLM architectures (Llama, Mistral, Falcon, and others)
    on Wikipedia $+$ code corpora.
\end{description}
\end{mdframed}

\section{The Big Picture}

Mechanistic interpretability seeks to decode the internal representations of large
language models---how concepts are organized, transformed, and composed across layers.
Dominant methods (probes, sparse autoencoders, dictionary learning) rest on a shared
linearity assumption: each concept occupies a direction in activation space, and its
presence is measured by a linear classifier or inner product.

This assumption is increasingly challenged by empirical evidence showing curved,
non-linearly distributed activation clusters. If concepts genuinely occupy
low-dimensional manifolds rather than directions, linear methods will systematically
miss or mischaracterize them. This paper takes the manifold hypothesis seriously
and develops tools to act on it, revealing organizational principles in LLMs
that linear baselines cannot detect.

\section{Why Existing Approaches Fall Short}

\textbf{Linear probes} measure concept presence as a dot product with a learned
direction, entirely missing curved or non-linearly distributed concept regions.

\textbf{Sparse autoencoders (SAEs)} decompose activations into sparse linear
combinations; each dictionary feature is a direction, and compositions are linear
sums. Non-linear manifold structure is encoded only incidentally.

\textbf{CKA and PCA-based comparisons} compare linear subspaces across layers
using canonical angles; they are insensitive to non-linear geometry and cannot
detect manifold-level structural differences.

\textbf{No cross-model alignment:} Existing methods lack a rigorous way to compare
concept organization across different architectures without explicit feature
matching, which assumes a known bijection between concepts---not generally available.

\section{Core Mechanism}

\textbf{NLMCD on LLM Activations.} For each concept category $c$ (syntactic
POS tag, dependency role, semantic entity type) and each layer $l$, the paper
collects token activation sets $\mathcal{H}_c^l$ and fits a low-dimensional
Riemannian manifold $\mathcal{M}_c^l$ using a neural network atlas.

\textbf{Concept-Based Alignment (CBA) Score.} CBA measures whether two manifolds
partition the same token set into the same neighborhoods, without requiring
explicit feature matching:
\[
\text{CBA}(\mathcal{M}_c^l, \mathcal{M}_c^{l'})
= \text{RandIndex}(G_c^l,\, G_c^{l'})
\]
where $G_c^l$ is a neighborhood graph on tokens $\mathcal{T}_c$ derived from
geodesic distances on $\mathcal{M}_c^l$. CBA is 1.0 when both manifolds partition
the token space identically and 0.5 at chance.

\section{Technical Formulation}

Let $\mathbf{h}_l^{(t)} \in \mathbb{R}^d$ be the activation of token $t$ at layer $l$.
The NLMCD atlas $\phi_c^l : \mathbb{R}^k \to \mathbb{R}^d$ (with $k \ll d$) is fitted by:
\[
\min_{\phi} \sum_{t \in \mathcal{T}_c}
  \|\mathbf{h}_l^{(t)} - \phi(z^{(t)})\|^2 + \lambda\,\mathrm{Reg}(\phi)
\]
where $z^{(t)} \in \mathbb{R}^k$ are latent coordinates learned jointly, and
$\mathrm{Reg}(\phi)$ is a smoothness regularizer.

The CBA between concept $c$ at layers $l$ and $l'$ uses the generalized Rand index:
\[
\text{CBA} = \frac{|\{(i,j): \text{near in }G_c^l \Leftrightarrow \text{near in }G_c^{l'}\}|}{|\mathcal{T}_c|^2}
\]
All-pairs CBA matrices across layers $\{1, \ldots, L\}$ reveal which layer intervals
share similar concept geometry.

\section{Learning or Inference Procedure}

\textbf{Manifold fitting (offline):}
\begin{enumerate}[itemsep=2pt,parsep=0pt]
  \item Run the LLM on a large corpus ($\approx$10M tokens).
  \item Extract activations $\{\mathbf{h}_l^{(t)}\}$ for all tokens and layers.
  \item For each concept category $c$ and layer $l$, fit NLMCD atlas $\phi_c^l$
    using stochastic gradient descent.
\end{enumerate}

\textbf{CBA computation (post-fitting):}
\begin{enumerate}[itemsep=2pt,parsep=0pt]
  \item Compute pairwise geodesic distances on each fitted manifold.
  \item Build neighborhood graphs $G_c^l$ and compute Rand index between
    all (layer, model) pairs.
  \item Visualize as CBA alignment matrices to identify block structures.
\end{enumerate}

\section{What the Guarantee Says}

NLMCD provides a consistent manifold estimate as sample size grows, under standard
regularity conditions on manifold geometry. CBA is a metric on the space of
neighborhood partitions. The main empirical claim---cross-architecture consistency
of block structures---is validated by replicating CBA patterns across five model
families including architectures trained on different data regimes.

\section{Experimental Findings}

\begin{itemize}[itemsep=2pt,parsep=0pt]
  \item \textbf{Block structures:} CBA matrices show consistent high-alignment blocks
    at intermediate layers ($\approx$layers 10--20 in a 32-layer model).
  \item \textbf{Syntax-to-semantics:} Syntactic concepts dominate CBA structure
    through $\approx$70\% of network depth; semantic concepts become prominent
    in the final 30\% of layers across all tested models.
  \item \textbf{Cross-architecture:} Block structures are consistent across Llama,
    Mistral, Falcon and two additional families at equivalent depth fractions.
  \item \textbf{CBA sensitivity:} CBA is 15--30\% more sensitive than PCA-based
    or CKA-based baselines, as measured by detection of synthetic perturbations.
  \item \textbf{Non-linearity:} Linear probes in place of NLMCD miss the block
    transitions entirely, confirming that non-linear geometry carries genuine
    structural information.
\end{itemize}

\section{Ablations and Interpretation}

\textbf{Manifold dimension $k$:} Values $k \in \{2,4,8\}$ all produce similar CBA
patterns; $k=4$ offers the best fidelity-cost trade-off.

\textbf{Corpus composition:} Code-heavy vs.\ text-heavy corpora show similar
syntactic block patterns but differ in semantic concept structure of late layers,
aligned with known training data differences.

\textbf{Tokenization:} Character-level tokens show less pronounced block structure
than subword tokens, suggesting the patterns are tightly coupled to the tokenization
scheme.

\section*{Reference}

Tido Specht, Elias Benedict Krey, Nils Neukirch, Nils Strodthoff.
\textbf{Beyond Linear Concepts: Discovering and Aligning Non-Linear Concept Manifolds
in Large Language Models}.
arXiv:2610.01821, October 2026.\\
\url{https://arxiv.org/abs/2610.01821}

\end{document}
